\section{Problema}
Su servicio de streaming de música le recomienda una lista de reproducción con $n$ canciones nuevas: $c_1, c_2, \ldots c_n$. Al escucharla, usted no necesariamente sigue el orden recomendado: las escucha en el orden que le da la gana. Para saber qué tanto difiere el orden en el que usted escuchó las canciones del orden propuesto, se cuenta la cantidad de pares invertidos: cuántas veces sucede que $c_j$ precede a $c_i$ para $i < j$.

Diseñe un algoritmo de dividir y conquistar que cuente el número de pares invertidos.
\begin{itemize}
  \item Entrada: Un arreglo \inline{canciones} de longitud \(n\) con enteros entre $1$ y $n$, donde cada entero $k$ representa la canción $c_k$ y la posición en el arreglo corresponde a la posición en la que realmente se escuchó.
  \item Salida: el número de pares invertidos.
\end{itemize}

E.g., el arreglo \inline{[3, 1, 4, 5, 2]} corresponde a la secuencia \(c_3, c_1, c_4, c_5, c_2\), que contiene cuatro pares invertidos: \((c_3, c_1)\), \((c_3, c_2)\), \((c_4, c_2)\) y \((c_5, c_2)\).

\section{Entregables}

\begin{enumerate}
  \item (0 puntos) Si lo desea, use el espacio restante en esta página para realizar una lluvia de ideas sobre su algoritmo. Puede escribir o dibujar cualquier cosa.
  \item (5 puntos) En la cuadrícula en el revés de la hoja, implemente su algoritmo en Python, Java, C, C++, JavaScript o TypeScript.
\end{enumerate}
